% Interpret the results, discuss limitations, and outline future improvements.

The experimental results reveal a nuanced performance landscape in
which solution quality, computational cost, and constraint
satisfaction trade off in ways that both confirm and extend the
findings reported in the orienteering literature.

\subsection{Interpretation of Results}

The most striking result is the parity between GRASP and Tabu
Search at the highest observed tour quality of 98.7, both visiting
11 landmarks within the eight-hour budget, yet with a wall-clock
time ratio of approximately 27:1 in favour of Tabu Search (0.37
seconds versus 10.18 seconds). This finding is consistent with the
comparative analysis of \textcite{Liang2002}, who observed that
Tabu Search and population-based heuristics achieve comparable
solution quality on orienteering benchmarks, and with
\textcite{Tang2005}, who demonstrated that adaptive memory
mechanisms in Tabu Search enable rapid convergence through
structured neighbourhood pruning. The successive-filtering
candidate strategy implemented in our Tabu Search --- ranking
landmarks by induced travel cost relative to interest score ---
appears to be the principal factor responsible for its
computational efficiency, effectively reducing the neighbourhood
evaluation burden without sacrificing solution quality.

Simulated Annealing with the Cauchy acceptance criterion (96.9,
11 landmarks) consistently outperforms the Boltzmann variant
(95.3, 11 landmarks), a result attributable to the heavier tails
of the Cauchy distribution, which sustain a higher probability of
accepting deteriorating moves at intermediate temperatures and
thereby reduce premature convergence. This echoes the theoretical
expectations discussed by \textcite{Gunawan2016}, who note that
acceptance-function design is a primary differentiator of
Simulated Annealing variants on constrained routing problems.

The Genetic Algorithm (68.3--90.0 across configurations) exhibits
greater variance than single-solution methods, which is consistent
with the sensitivity of population-based approaches to crossover
operator design on time-windowed instances. The order crossover
and tailored crossover implementations produced different fitness
landscapes, and the relatively modest performance of the penalty
fitness function relative to the feasibility-oriented variant
suggests that soft-constraint handling significantly affects
search trajectory for this problem class.

The most academically interesting result is the underperformance
of the CPLEX exact solver (61.8--88.2, 7--10 landmarks) relative
to GRASP and Tabu Search. Three explanations are plausible. First,
CPLEX required 30.2 seconds to terminate, suggesting that the
branch-and-bound search tree was not fully explored within the
available computation budget, and the best integer solution found
at termination is consequently a feasible but sub-optimal
incumbent. Second, the Mixed Integer Linear Programming
formulation introduces symmetry constraints and big-$M$
linearisations for time-window propagation that tighten the
feasible region relative to the heuristic representations. Third,
the haversine-based travel matrix, while consistent across all
solvers, produces a non-integer cost structure that may slow LP
relaxation bounds. \textcite{Righini2009} observed similar
behaviour and addressed it through decremental state-space
relaxation; incorporating such techniques would likely close the
gap between CPLEX and the leading metaheuristics.

\subsection{Limitations}

Several limitations constrain the generalisability of these
results. First, travel times are computed from straight-line
haversine distances at a fixed average urban speed of 25 km/h.
Real road distances in Algiers, which the backend routing service
retrieves from MAPBOX API, are systematically longer and temporally
variable due to traffic. The algorithms were therefore optimising
a surrogate objective that diverges from the true tourist
experience. Second, landmark interest scores were assigned through
a weighted combination of public ratings and expert judgement, a
process that introduces subjectivity and limits
reproducibility. Third, the dataset is confined to 56 landmarks
in the Algiers wilaya, a scale at which all tested algorithms
remain tractable; results may not extrapolate to larger
metropolitan datasets. Finally, the time-window representation
assumes deterministic opening hours, whereas real tourist sites
may close unexpectedly or have dynamic crowd-dependent effective
visit durations.

\subsection{Future Work}

Several directions offer meaningful improvements. Within the
classical metaheuristic paradigm, Adaptive Large Neighbourhood
Search (ALNS) represents the current state of the art for
time-constrained orienteering \parencite{Gunawan2016}, with its
self-tuning destroy-and-repair operators addressing the principal
weakness of both GRASP (construction diversity) and Tabu Search
(neighbourhood breadth). Iterated Local Search with the MaxShift
data structure of \textcite{Vansteenwegen2009} would provide an
efficient feasibility oracle that reduces the per-iteration cost
of neighbourhood evaluation from $O(n)$ to $O(1)$.

At the frontier of the field, deep reinforcement learning (DRL)
approaches have demonstrated near-optimal performance on routing
problems of this scale. \textcite{Kool2019} showed that an
Attention Model trained with the REINFORCE policy gradient
algorithm learns generalisable construction heuristics for TSP and
related problems; \textcite{Gama2020} extended this framework
directly to the OPTW. Hybrid DRL-ALNS architectures, in which a
learned policy selects neighbourhood operators dynamically
\parencite{Reijnen2024}, represent the most promising direction
for achieving both solution quality and sub-second inference
times. A further extension of direct practical relevance is the
Stochastic Orienteering Problem, in which travel times are drawn
from a distribution reflecting real traffic conditions; this
formulation is underexplored in the literature
\parencite{Vansteenwegen2011} and maps naturally onto the
Algiers urban context. Finally, extending the single-day
formulation to the Team Orienteering Problem with Time Windows
\parencite{Gunawan2016} would enable multi-day, multi-guide
itinerary planning, substantially increasing the practical utility
of the system.
